% Task I outline: Life-cycle sustainability of AI inference measured per token
% Working hypothesis: tokens may be a useful activity unit, but not necessarily
% a sufficient functional unit because they do not directly measure user value.

\documentclass[11pt,a4paper]{article}

\usepackage[margin=2.5cm]{geometry}
\usepackage{amsmath}
\usepackage{booktabs}
\usepackage{enumitem}
\usepackage{hyperref}

\title{Is a Token a Meaningful Functional Unit for the\\
Life-Cycle Assessment of AI Inference?}
\author{Author 1 \and Author 2}
\date{\today}

\begin{document}

\maketitle

% Replace this with a short summary once the scope and data sources are fixed.
\begin{abstract}
\textit{TODO: State the case, the proposed functional unit, the main system
boundary, the expected contribution, and the main uncertainty in 100--150
words.}
\end{abstract}

\tableofcontents

\section{Introduction}

\subsection{Background and motivation}
% Explain why this case matters in the context of the course case study.
% Suggested points:
\begin{itemize}[leftmargin=*]
    \item AI services depend on large-scale computation, electricity, cooling,
    networking, and semiconductor hardware.
    \item Digitalisation and AI may increase demand for microchips and data-centre
    capacity, linking the case to the environmental impacts of microfabrication.
    \item The environmental burden of AI is often reported per request or per
    token, but the basis for this choice is not always transparent or comparable.
    \item A functional unit should represent the function delivered by the system,
    not only an easily counted activity.
\end{itemize}

\subsection{Connection to microfabrication}
% Decide explicitly how close the study should be to the course's microfabrication
% case. Keep one of the following framings and justify the choice.
\begin{itemize}[leftmargin=*]
    \item \textbf{Primary framing:} AI inference is assessed as a digital service
    enabled by semiconductor manufacturing; chip and data-centre impacts are
    included where data permits.
    \item \textbf{Narrow framing:} the study assesses only operational inference
    energy, while semiconductor manufacturing and model training are discussed as
    excluded upstream processes and uncertainty sources.
\end{itemize}

\subsection{Problem statement}
% Draft the problem statement as one paragraph. It should identify the decision
% or knowledge gap, not merely describe the topic.
\textit{Working draft:} It is unclear whether ``one token'' is a sufficiently
meaningful and reproducible functional unit for comparing the sustainability of
AI language-model inference. Token counts are measurable, but they may not
capture differences in model capability, response quality, input context,
hardware, electricity mix, or the useful service delivered to a user.

\subsection{Aim and research questions}
\textbf{Overall aim:} To evaluate whether a token-based functional unit can
support a transparent and decision-relevant sustainability assessment of AI
language-model inference.

\begin{enumerate}[leftmargin=*]
    \item What is the object of assessment and what service does an AI inference
    system provide?
    \item How should the functional unit be defined so that input tokens, output
    tokens, model quality, and workload are represented?
    \item What environmental, economic, and social factors dominate the result?
    \item How sensitive are the results to model size, hardware utilisation,
    response length, data-centre electricity mix, and system boundaries?
    \item Does a token-based unit provide a more useful comparison than an
    alternative unit such as one completed task or one accepted answer at a
    specified quality level?
\end{enumerate}

\section{Task I -- Problem formulation}

\subsection{Object of assessment}
% Define the system in operational terms. Avoid using ``AI'' as the object by
% itself; name the service and the life-cycle stage being assessed.
\begin{itemize}[leftmargin=*]
    \item \textbf{Object:} an AI language-model inference service that receives a
    user prompt and returns generated text.
    \item \textbf{Service function:} transformation of an input prompt into a
    useful generated response under specified quality and latency conditions.
    \item \textbf{Reference flow:} the electricity, cooling, hardware use,
    networking, and other inputs required to deliver the defined service.
    \item \textbf{Unit process:} one inference event, or a batch of inference
    events, at a named model and hardware configuration.
\end{itemize}

\subsection{Functional unit}

\subsubsection{Candidate definitions}
% The primary proposal should be measurable and reproducible. Compare it with
% alternatives before selecting the final unit.
\begin{table}[htbp]
    \centering
    \begin{tabular}{p{0.27\textwidth} p{0.34\textwidth} p{0.29\textwidth}}
        \toprule
        Candidate & Strength & Limitation \\
        \midrule
        1 million output tokens & Simple, measurable, and easy to scale &
        Ignores input length, quality, model capability, and usefulness \\
        1 million total tokens (input + output) & Includes both sides of the
        inference workload & Input and output tokens may have different energy
        requirements; still not a direct measure of value \\
        One standardised prompt--response pair & Represents a concrete workload
        and enables model comparison & Results depend strongly on the selected
        prompt, response length, and quality \\
        One completed task at a defined quality level & Closest to the service
        delivered to a user & Quality is harder to measure and may require human
        evaluation \\
        \bottomrule
    \end{tabular}
    \caption{Functional-unit candidates to evaluate before selecting the main unit.}
    \label{tab:functional-unit-candidates}
\end{table}

\subsubsection{Proposed working functional unit}
% Select and justify one unit after the literature/data review. A good starting
% point is deliberately specific rather than simply ``one token''.
\textbf{Working proposal:} one million tokens processed by a specified AI
inference workload, defined by an input-to-output token ratio, model class,
quality target, hardware configuration, and data-centre electricity mix.

\begin{equation}
    FU = 10^{6}\ \text{tokens processed}
    \quad\text{with a specified input/output ratio and service quality.}
\end{equation}

% Explain why the unit is suitable, what it leaves out, and how the result will
% be stress-tested against ``one completed task at defined quality''.
\textit{Decision to make:} The final report should state whether tokens are being
used as (i) the functional unit, (ii) an intermediate activity unit, or (iii) a
scenario variable that must be combined with a quality-based functional unit.

\subsection{Quantification of the function}
% Describe how the function will be measured. Do not claim that token count alone
% measures usefulness.
\begin{itemize}[leftmargin=*]
    \item Input tokens and output tokens per inference event.
    \item Model identity, parameter/model class, quantisation, and serving mode.
    \item Hardware type, accelerator utilisation, batch size, and assumed lifetime.
    \item Response quality or task success criterion, if a quality-based unit is
    selected.
    \item Latency/service availability assumptions, if they influence the function.
\end{itemize}

\subsection{System boundary}

\subsubsection{Included processes}
% Mark each item as included, excluded, or sensitivity-only once data sources are
% known.
\begin{itemize}[leftmargin=*]
    \item Electricity consumed during inference.
    \item Cooling and other data-centre overheads, represented by PUE or a similar
    factor.
    \item Manufacturing and replacement of servers/accelerators, allocated over
    useful lifetime and workload where data permits.
    \item Networking and storage associated with the service, if relevant.
    \item Upstream electricity generation and regional electricity mix.
\end{itemize}

\subsubsection{Excluded or separate scenario processes}
\begin{itemize}[leftmargin=*]
    \item Model development and training, unless an amortisation scenario is
    added.
    \item Construction of data-centre buildings and semiconductor fabrication,
    unless reliable data can support their inclusion.
    \item End-user device energy, unless the service boundary is extended to the
    full user interaction.
    \item Rebound effects and changes in demand, treated as a limitation or
    separate systemic scenario rather than hidden in the operational result.
\end{itemize}

\subsection{Reference flow and calculation structure}
% Keep the calculation transparent and dimensional. Replace symbols with data in
% Task II.
\begin{align}
    E_{FU} &= E_{\mathrm{inference}} + E_{\mathrm{cooling}} +
    E_{\mathrm{network}} + E_{\mathrm{allocated\ hardware}}, \\
    CF_{FU} &= E_{FU} \times EF_{\mathrm{electricity}},
\end{align}
where $E_{FU}$ is energy per functional unit and $CF_{FU}$ is the resulting
carbon footprint per functional unit. The allocation method and emission factor
must be reported explicitly.

\subsection{Three-pillar motivation}

\subsubsection{Environmental motivation (Planet)}
% Energy use, greenhouse-gas emissions, water use, critical materials, e-waste,
% and upstream semiconductor impacts.

\subsubsection{Economic motivation (Prosperity)}
% Cost of electricity, hardware utilisation, infrastructure investment, and the
% economic value of more efficient models or shorter responses.

\subsubsection{Social motivation (People)}
% Access to AI services, labour and supply-chain conditions, privacy, digital
% inequality, and the societal consequences of scaling AI services.

\subsubsection{Potential trade-offs}
% Examples: larger models may improve quality but increase energy demand; longer
% responses may improve task completion but increase token use; local deployment
% may reduce latency but reduce utilisation; efficiency measures may increase
% hardware demand or critical-material use.

\subsection{Assumptions and limitations}
% State assumptions before collecting data. Use a table if the list becomes long.
\begin{itemize}[leftmargin=*]
    \item Tokenisation is consistent within the selected model and language.
    \item The reported energy is allocated consistently between inference events.
    \item Hardware lifetime and utilisation are known or varied in sensitivity
    analysis.
    \item Electricity emission factors represent the selected geography and time
    period.
    \item Model quality is either held constant, measured explicitly, or clearly
    acknowledged as missing from a token-only unit.
\end{itemize}

\subsection{Uncertainty plan}
% This section is required by the evaluation criteria and should be designed in
% Task I, even though the numerical analysis belongs in Task II.
\begin{itemize}[leftmargin=*]
    \item Parameter uncertainty: power draw, PUE, utilisation, lifetime, and
    electricity emission factor.
    \item Scenario uncertainty: training included vs excluded; hardware allocation;
    regional electricity mix; model size; input/output ratio.
    \item Structural uncertainty: token-based unit versus quality-based service
    unit; attributional versus broader systemic perspective.
    \item Data quality: vendor estimates, academic measurements, ESG reports, and
    secondary databases should be ranked by reliability and comparability.
\end{itemize}

\section{Planned methods for Task II}
% Keep this section as a bridge. It can be expanded after Task I is approved.
\begin{enumerate}[leftmargin=*]
    \item Define a standard workload and collect energy or power data.
    \item Calculate energy and carbon footprint per candidate functional unit.
    \item Perform sensitivity analysis for the dominant assumptions.
    \item Compare the result with other digital services or relevant consumables,
    while documenting differences in function and quality.
    \item Evaluate economic and social indicators alongside the environmental result.
\end{enumerate}

\section{Planned systemic assessment for Task III}
% Placeholder for the later integrated assessment.
\begin{itemize}[leftmargin=*]
    \item Relate the result to planetary boundaries or another global benchmark.
    \item Discuss circular-economy options for hardware, heat, water, and
    electricity systems.
    \item Assess trade-offs between environmental, economic, and social outcomes.
    \item Explain what the chosen functional unit makes visible and what it hides.
    \item Integrate Task I and Task II as appendices in the final submission.
\end{itemize}

\section{Expected contribution}
% End Task I by stating what the study will deliver.
\textit{TODO: Explain whether the main contribution is a numerical estimate, a
comparison of functional units, a framework for transparent reporting, or a
combination of these.}

\section{References}
% Add sources from DTU Learn and the literature review. Use a .bib file when the
% bibliography is established.
% \bibliographystyle{plain}
% \bibliography{references}

\appendix
\section{Appendix: Decision log for scope and functional unit}
% Record decisions, rejected alternatives, and the reason for each decision.

\section{Appendix: Data collection template}
% Suggested fields: source, date, model, hardware, workload, input tokens,
% output tokens, power, runtime, PUE, geography, emission factor, uncertainty.

\end{document}
